\documentclass[10pt,a4paper]{amsart}
\usepackage[margin=19mm]{geometry}
\usepackage{amsmath,amssymb,mathtools}
\usepackage[scr=rsfs,cal=euler]{mathalfa}
\usepackage{xfrac}
\usepackage[unicode,hidelinks,bookmarks=false]{hyperref}
\newcommand{\ie}{i.e.\ }
\newcommand{\cf}{cf.\ }
\newcommand{\resp}{respectively\ }
\newcommand{\Subsection}{\subsection}
\providecommand{\nobreakdash}{\nobreak\hbox{-}}
\setlength{\emergencystretch}{2em}
\multlinegap0pt
\usepackage{bm}
\usepackage{bbm}
\usepackage{dsfont}
\usepackage{xspace}
\usepackage{cancel}
\usepackage{mathtools}
\usepackage{amsthm}
\makeatletter
\@ifundefined{theorem}{\newtheorem{theorem}{Theorem}}{}
\@ifundefined{lemma}{\newtheorem{lemma}[theorem]{Lemma}}{}
\makeatother
\title[Analytic structure in spaces of Lipschitz functions]{Analytic structure in spaces of Lipschitz functions}
\author{Stephen Deterding}
\date{Source: arXiv:2501.19181v2 (CC BY 4.0)}
\begin{document}
\maketitle

\noindent\emph{Source and licence.} Analytic structure in spaces of Lipschitz functions. Version arXiv:2501.19181v2 (CC BY 4.0), distributed under the Creative Commons Attribution 4.0 International licence, as recorded in the arXiv licence metadata for this exact version and shown on the abstract page https://arxiv.org/abs/2501.19181v2. Full rights notice: licenses/arxiv-2501.19181-CC-BY-4.0.txt.\par\smallskip

\noindent\textit{Editorial scope.} Counted unit: the Theorem whose source label is \texttt{c then b} in the article (source0.tex lines 304--326, source label \texttt{c then b}), delivered together with the article's own complete original proof of it (source0.tex lines 329--401, one \texttt{proof} environment of 73 lines). No mathematics is written, repaired or re-derived: statement and proof are the article's own text line for line. Required context retained uncounted: Lord (lines 285--296). Every definition, display and label of those lines is retained unchanged. Omitted from this package and not redistributed: 1--284, 297--303, 327--328, 402--751. Nothing inside a retained excerpt is omitted or altered: every retained body is delivered exactly as the article's lines are written, so no omission receipt of any kind accompanies this package. Citations: the retained excerpts carry no citation command at all, so no bibliography entry of the article is transcribed and no reference is redistributed. Reference adaptation: none was needed and none was made; every reference inside the delivered unit resolves inside it, so no unresolved reference or citation is produced. Printed slips kept literal: no printed word, symbol, hypothesis or number of the counted statement or of its proof was altered. Layout and class: the article's own class is \texttt{article} with packages including latexsym, amssymb, amsmath, amsfonts, pictex, enumerate, amsthm, dsfont, mathrsfs, tikz, scalerel, stackengine, srcltx; the wrapper is the lane's pinned \texttt{amsart} editorial wrapper at 10pt on A4, which restates the theorem-like environments the retained text needs and adds the article's own macro definitions that the retained text needs (none), with extra wrapper packages (bm, bbm, dsfont, xspace, cancel, mathtools). The delivered numbering is the unit's own. Assets: no retained span contains a figure environment, a graphics-inclusion command, a PDF-page inclusion, a table or a TikZ picture; no image file is copied into this package, the package declares no asset dependency and no image file is redistributed with it. Licence: version arXiv:2501.19181v2 of this article is distributed under the Creative Commons Attribution 4.0 International licence, as recorded in the arXiv licence metadata for this exact version and shown on the abstract page https://arxiv.org/abs/2501.19181v2. A full rights notice is delivered at licenses/arxiv-2501.19181-CC-BY-4.0.txt. Characters: every retained excerpt is pure ASCII, so the delivered engine, XeLaTeX with its default Latin Modern text font, reports no missing character.
\begin{lemma}
\label{Lord}
     Let $\Gamma$ be a piecewise analytic curve that encloses a region $U$ and is free of outward pointing cusps, and let $0<\alpha<1$. Then there is a constant $C$ which only depends on the curve $\Gamma$ such that

\begin{equation*} 
\left| \int_{\Gamma} f(z) dz\right| \leq C M_*^{1+\alpha}(K) ||f||_{\text{Lip}_\alpha(K)}'.
\end{equation*}

\noindent whenever $f \in $ lip$_\alpha(\mathbb C)$
 is analytic on $U \setminus K$, where $K$ is a compact subset of $U$.

\end{lemma}

\begin{theorem}
\label{c then b}
 Let $\phi(r)$ be a positive non-decreasing function, let $U$ be an open subset of the complex plane with $x \in U$, and let $Y$ denote the closure of $U$. Suppose

\begin{align*}
    \sum_{n=1}^{\infty} 2^n \phi(2^{-n}) M^{1+\alpha}_*(A_n(x) \setminus U) < \infty.
\end{align*}


\noindent Then there exists a Borel regular measure $\mu$ on $Y \times Y$ such that for all functions $f \in A_\alpha(U)$,

\begin{align*}
    f(x) = \int \frac{f(z)-f(w)}{|z-w|^{\alpha}} d\mu(z,w)
\end{align*}


\noindent and 

\begin{align*}
    \int \frac{d\mu(z,w)}{\phi(|z-x|)} < \infty.
\end{align*}
    
\end{theorem}

\begin{proof}
    Without loss of generality, we may assume that $x=0$. Let $A_n = A_n(0)$ and let $T_n$ be a linear functional on $A_\alpha(U \cap A_n)$ defined by

\begin{align*}
    T_n(f) = -\int_{A_n} \frac{f(z)}{z} dz
\end{align*}

\noindent and let $f \in A_\alpha(U)$. It follows from Lemma \ref{Lord} that 

\begin{align*}
    \left| \int_{A_n} \frac{f(z)}{z} dz \right| \leq C 2^n M_*^{1+\alpha}(A_n \setminus U) ||f||_{\text{Lip}_{\alpha}(A_n \setminus U)}'
\end{align*}


\noindent and thus $T_n$ is a bounded linear functional on $A_{\alpha}(U \cap A_n)$. Let $Y_n$ denote the closure of $U \cap A_n$. It follows from De Leeuw's representation of lip$\alpha^*$ that there is a Borel-regular measure $\nu_n$ on $Y_n \times Y_n$ such that

\begin{align*}
    T_n(f) = \int \frac{f(z)-f(w)}{|z-w|^{\alpha}} d\nu_n(z,w)
\end{align*}


\noindent and 

\begin{align*}
    \int |\nu_n(z,w)| \leq C 2^n M_*^{1+\alpha}(A_n \setminus U).
\end{align*}

\noindent Let $\displaystyle \nu = \frac{1}{2\pi i} \sum_{n=1}^{\infty} \nu_n$ and let $T$ be the linear functional on $A_\alpha(U \cap A_0)$ defined by 

\begin{align*}
    T(f) = \frac{1}{2\pi i}\int_{|z|=\frac{1}{2}} \frac{f(z)}{z} dz.
\end{align*}

\noindent Then $|T(f)| \leq 2||f||_{\text{Lip}_{\alpha}}(U \cap A_0)$ and thus $T$ is bounded. Hence, there is a Borel regular measure $\lambda$ on $Y_0 \times Y_0$ such that 

\begin{align*}
    T(f) = \int \frac{f(z)-f(w)}{|z-w|^{\alpha}} d\lambda(z,w)
\end{align*}

\noindent and 

\begin{align*}
    \int |\lambda(z,w)| \leq C.
\end{align*}


 Since $f$ is analytic in a neighborhood of $0$, there exists a positive integer $N$ such that $f$ is analytic on the disk $\{z:|z| \leq 2^{-N}\}$. Hence, by the Cauchy integral formula, there exists a constant $N>0$ such that $\displaystyle f(0) = \frac{1}{2\pi i}\int_{|z|=2^{-N}} \frac{f(z)}{z}dz$. Thus,

\begin{align*}
f(0) &= \frac{1}{2\pi i}\int_{|z|=2^{-N}} \frac{f(z)}{z}dz\\
 &= \frac{1}{2\pi i}  \int_{|z|=\frac{1}{2}} \frac{f(z)}{z} dz - \frac{1}{2\pi i} \sum_{n=1}^{\infty} \int_{\partial A_n} \frac{f(z)}{z} dz\\
    &= \int \frac{f(z)-f(w)}{|z-w|^{\alpha}} d\lambda(z,w) + \int \frac{f(z)-f(w)}{|z-w|^{\alpha}} \frac{1}{2\pi i} \sum_{n=1}^{\infty} d\nu_n(z,w)\\
     &= \int \frac{f(z)-f(w)}{|z-w|^{\alpha}} d\lambda(z,w) + \int \frac{f(z)-f(w)}{|z-w|^{\alpha}} d\nu.
\end{align*}

\noindent Let $\mu = \lambda + \nu$. Then 

\begin{align*}
    f(0) = \int \frac{f(z) - f(w)}{|z-w|^{\alpha}} d \mu
\end{align*}

\noindent as desired. Lastly, since $\nu_n$ has support on $Y_n \times Y_n$ it follows that

\begin{align*}
    \int \frac{|d\mu(z,w)|}{\phi(|z|)} &\leq \frac{1}{2\pi} \int \phi(|z|)^{-1} \left| \sum_{n=1}^{\infty} \nu_n(z,w) \right| + \int \phi(|z|)^{-1} |\lambda|\\
   % &\leq\frac{1}{2\pi}  \int \sum_{n=1}^{\infty} \nu_n(z,w) + C\\
    &\leq\frac{1}{2\pi} \int \sum_{n=1}^{\infty} \phi(|2^{-n}|)^{-1} |\nu_n(z,w)| + C\\
    &\leq C \sum_{n=1}^{\infty}  \phi(|2^{-n}|)^{-1} 2^n M_*^{1+\alpha}(A_n \setminus U) < \infty,
\end{align*}

\noindent thus completing the proof.
    
\end{proof}

\end{document}
